\section{Procedimiento y cobertura}\label{sec:comparacion:procedimiento}\label{sec:hallazgos:comparacion}

El catálogo se aplica a las ocho Datasheets en tres pasos.
Primero, cada una de las 32 preguntas se instancia sobre cada Datasheet con los hechos que esta registra, cada dato que se recogió con su fuente al construirla, esté o no en un campo.
Después, las preguntas de clase P, las que resuelve una operación de PRIME4API, se calculan a mano con la definición de la operación y se contrastan con la respuesta del motor.
Por último, cada pareja de pregunta y API recibe una de cuatro salidas.

La primera salida es la respuesta con su cifra y con la fuente que la sostiene.
La segunda es la respuesta parcial, cuando solo sale una parte de la pregunta, y la tercera es «no», con el dato que falta.
La cuarta, «fuera de los campos», solo existe en la clase X, la de las preguntas que el instrumental no responde.
Se anota cuando los hechos de la API contienen el dato aunque ningún campo de la Datasheet lo recoja, por ejemplo en el \textit{provenance} o en la matriz de \textit{features}.

El recuento acepta «no publicado» como respuesta cuando la pregunta lo admite.
Así se responde PI-A02, la velocidad máxima de envío, en Postmark.
La respuesta es que el proveedor no publica el \textit{rate}, junto con las 5.000 peticiones por segundo que la Datasheet toma de la cabecera \texttt{ratelimit-limit} leída en un sondeo, una prueba propia contra la API.
Ese valor está en ámbar (\semambar{}), el color que el semáforo de la Datasheet da a un valor interpretado.
En las preguntas de comparación de la categoría C, una API cuenta como respondida cuando aporta su lado de la comparación.

Un mes se toma como 30 días (2.592.000 segundos).
Cada respuesta vale a la fecha de su Datasheet, dentro de la ventana de toma de datos del 2 de agosto al 21 de septiembre de 2026 (Sección~\ref{sec:propuesta:ventana}).
Entre capturas sucesivas de una misma página cambió la redacción, pero ningún valor que la Datasheet modele.

Los términos de consumo de una API (el precio y la información de cómo, cuándo, cuánto y dónde se consume) están repartidos entre el panel de la cuenta y varias páginas del proveedor, como la de precios, la documentación de la API, el centro de ayuda o el blog.
Cuando la página de precios y otra de esas fuentes no coinciden, la comparación usa el valor de la página de precios, que es el que modela la Datasheet: las demás fuentes señalan la incoherencia, pero no corrigen el dato (Sección~\ref{sec:propuesta:semaforo}).

El catálogo de Datasheets registra 82 divergencias entre fuentes, cada una un dato que dos fuentes dan distinto, y en 37 la página de precios es uno de los dos lados.
En siete de esas 37 la página se contradice a sí misma, y la regla de la página de precios no ordena nada.
En once de las 30 restantes, esa regla cambió la fuente que ganaba antes de adoptarla o cerró una divergencia que seguía abierta.
Solo dos de esas once cambiaron la Datasheet: un valor, el sobrecoste de Essentials 50K en SendGrid, y un color del semáforo, el de la cuota diaria del plan gratuito de Mailgun.
Las otras nueve solo cambian la procedencia anotada, sin tocar ningún campo.

Algunas preguntas presuponen un plan con su cuota, y no todos los proveedores lo tienen.
PI-B02 pregunta cuántos correos compra un presupuesto mensual en un plan, y Mailchimp Transactional no tiene planes.
En esos casos la pregunta se responde sobre lo que la Datasheet sí declara, como una cuenta de prueba o una tarifa sin suscripción, y la celda se rotula entre corchetes con el nombre que le da el proveedor.
En Mailchimp es \texttt{[demo]}, como llama el proveedor a su cuenta de prueba.
En Amazon SES es \texttt{[sandbox]}, su entorno de pruebas y el único con cuota publicada, o \texttt{[à la carte]}, su tarifa sin suscripción, que cobra solo por correo enviado.
El corchete no lleva semáforo, porque no mide la fiabilidad del dato.

\paragraph{Ejecución contra PRIME4API.}
Las preguntas P se ejecutaron en 128 llamadas al motor, todas con su petición y su respuesta archivadas\footnote{En \texttt{05\_salida/\allowbreak \_cap5/\allowbreak prime4api\_ejecucion*/} del repositorio del trabajo.}.
Fueron 45 llamadas para Postmark, Resend y SendGrid, 15 para MailerSend y 68 para las 40 instancias de Brevo, Amazon SES, Mailgun y Mailchimp.
En las cuatro primeras APIs el motor contestó 58 de las 60 llamadas, y las dos que fallan piden PI-B01 con todos los planes y tropiezan con el plan Enterprise de Resend y con el de SendGrid.
En las otras cuatro contestó las 17 llamadas de Mailgun y las 17 de Amazon SES, 16 de 17 en Brevo y 3 de 17 en Mailchimp.
De sus 40 instancias resuelve 30, las diez de Brevo, las diez de Amazon SES y las diez de Mailgun.

Las respuestas del motor coinciden con el cálculo a mano salvo por las tres convenciones de conteo de la Sección~\ref{sec:limites:prime4api}.
Una de ellas abre la ventana siguiente en el instante final del intervalo.
Por eso, en PI-B02, el motor da 3.100 correos en el plan gratuito de Mailgun, cien más que su cuota mensual: suma 30 ventanas diarias de 100 correos y la que abre en $t = 30$ días, que ya es del mes siguiente.

Entre esas otras cuatro, el motor devuelve errores en Brevo y en Mailchimp.
La única llamada de Brevo que no nombra los planes alcanza Enterprise, cuyo precio es \texttt{custom}, y devuelve HTTP 400, el caso de un precio que existe y no es público que recoge P-05 (Sección~\ref{subsec:evolucion:p05}).
Las otras dieciséis nombran los planes y lo dejan fuera.
En Mailchimp, catorce llamadas devuelven HTTP 500 por dos campos vacíos que el esquema admite: el tope de destinatarios y el periodo de la franquicia de 500 correos de la cuenta de prueba (Sección~\ref{sec:esquema:contrato}).
Las tres que contesta solo leen su cuota o su caudal, y por eso sus diez respuestas de clase P están calculadas a mano, sin confirmación del motor, y la celda lo indica en cada caso.

\paragraph{Cobertura.}
De las 25 preguntas de clase D, P y H, las que el catálogo declara respondibles con una Datasheet, PRIME4API o HARVEY4API, Resend responde 22 con los hechos de su Datasheet.
Amazon SES y MailerSend responden 21, SendGrid, Mailgun, Mailchimp y Brevo 20 y Postmark 18.
Dieciséis preguntas se responden en las ocho APIs a la vez.
La Tabla~\ref{tab:cobertura-pi} recoge el recuento completo por API.

\begin{table}[tbp]
\centering
\scriptsize
\setlength{\tabcolsep}{2pt}
\begin{tabular}{|p{4.5cm}|c|c|c|c|c|c|c|c|c|}
\hline
\textbf{Salida} & \rotatebox{90}{\textbf{Postmark}} & \rotatebox{90}{\textbf{Resend}} &
\rotatebox{90}{\textbf{SendGrid}} & \rotatebox{90}{\textbf{MailerSend}} &
\rotatebox{90}{\textbf{Brevo}} & \rotatebox{90}{\textbf{SES}} &
\rotatebox{90}{\textbf{Mailgun}} & \rotatebox{90}{\textbf{Mailchimp}} &
\rotatebox{90}{\textbf{Las ocho}} \\
\hline
Respondible con los hechos & 18 & 22 & 20 & 21 & 20 & 21 & 20 & 20 & 16 \\
\hline
\quad de ellas, con un dato del \textit{provenance} y no de un campo del esquema & 4 & 3 & 2 & 1 & 5 & 6 & 1 & 2 & 1 \\
\hline
Parcial & 4 & 2 & 4 & 4 & 4 & 4 & 5 & 3 & 2 \\
\hline
Respondible en el catálogo, sin respuesta en los hechos & 3 & 1 & 1 & 0 & 1 & 0 & 0 & 2 & 0 \\
\hline
L, respondida leyendo la matriz de \textit{features} de la Datasheet & 2 & 2 & 2 & 2 & 2 & 2 & 2 & 2 & 2 \\
\hline
X sin respuesta en ningún sitio & 2 & 2 & 2 & 2 & 2 & 3 & 1 & 2 & 1 \\
\hline
X con respuesta en los hechos (fuera de los campos) & 3 & 3 & 3 & 3 & 3 & 2 & 4 & 3 & 2 \\
\hline
Total & 32 & 32 & 32 & 32 & 32 & 32 & 32 & 32 & --- \\
\hline
\end{tabular}
\caption[Cobertura del catálogo de preguntas sobre las ocho APIs]{Cobertura del catálogo de preguntas sobre las ocho APIs. El catálogo declara 25 respondibles con una Datasheet, PRIME4API o HARVEY4API (clases D, P y H), 2 que se responden leyendo (L) y 5 que el instrumental no responde (X), y ese reparto no cambia de una API a otra. La columna «Las ocho» cuenta las preguntas con la misma salida en las ocho APIs y no suma 32. La lista de preguntas de cada celda está en \texttt{05\_salida/\allowbreak \_cap5/\allowbreak analisis.md} del repositorio del trabajo.}
\label{tab:cobertura-pi}
\end{table}

De las dieciséis preguntas comunes, cinco piden cuota, precio o si el \textit{rate} está publicado (PI-A02, B03, C02, C03 y G03), y siete necesitan el valor del \textit{rate} (PI-A06, A07, B05, C04, D01, D02 y D03).
Las cuatro restantes piden el precio de un volumen (PI-B01, B02, C01 y G02).
En Brevo esas cuatro se responden con los volúmenes contratados de Starter, que van de 5.000 correos por 9 USD a 100.000 por 82 USD y constan en el \textit{provenance}, no en un campo (P-01, Apartado~\ref{subsec:evolucion:p01}).
El selector de volumen de su página de precios solo deja sin precio la última posición, «More», por encima de 100.000 correos en Starter, de 500.000 en Standard y de diez millones en Professional.

Los cuatro parciales de Brevo salen de lo que el proveedor no publica.
PI-A01, la cuota y lo que ocurre al superarla, queda en parcial porque sus tres planes de pago no lo documentan y la página de precios niega que exista precio de exceso \citep{brevopricing}.
Los otros tres salen de los tres tramos de caudal de su API, General, Advanced y Extended.
Professional puede acogerse a dos y Enterprise a los tres, y ninguna fuente dice cuál se les aplica.

\paragraph{Un caso resuelto: PI-B02.}
PI-B02, cuántos correos se envían como máximo con 50 USD al mes, recorre el procedimiento completo.
Resend llega a 116.000 correos con el segundo volumen de Pro (100.000 por 35 USD) más 16 bloques de 1.000 a 0,90 USD.
SendGrid llega a 116.722 con Essentials 100K, de 34,95 USD, más 15,05 USD de sobrecoste a 0,0009 USD por correo.
MailerSend con Starter y Mailgun con Foundation llegan los dos a 61.538: 50.000 correos por 35 USD y 15 USD a 1,30 USD el millar.

Las respuestas de Postmark y de Brevo llevan una salvedad cada una.
Postmark se queda en 36.666 con Platform (18 USD de plan y 32 USD a 1,20 USD el millar), sobre una cuota marcada en ámbar (\semambar{}).
Brevo llega a 40.000 correos con el volumen contratado de Starter que cuesta 46 USD, un precio que consta en el \textit{provenance} y no en un campo (P-01).

En las otras dos APIs el presupuesto no fija la respuesta.
En Amazon SES los 50 USD compran 500.000 destinatarios a la tarifa sin suscripción, pero el entorno de pruebas solo deja enviar 200 correos al día, 6.000 al mes, así que la respuesta la fija la capacidad y no el precio.
En Mailchimp compran 50.000 correos de cuota, pero la cuenta de pago no tiene caudal publicado, y con los 25 correos por hora de la cuenta de prueba el mes da 18.000, rotulados \texttt{[demo]}.

PRIME4API da 83.343 correos en Resend y 83.115 en SendGrid, frente a los 116.000 y 116.722 del cálculo a mano, porque de cada plan solo ve el primer volumen contratado, el de 50.000 correos (Sección~\ref{sec:limites:prime4api}).
Los 100.000 correos de Pro en Resend los vende la página de precios, y ACM2YAML v0.3 no tiene campo para guardarlos (P-01, los volúmenes contratados de un plan, Sección~\ref{subsec:evolucion:p01}).

\paragraph{Lectura de las columnas.}
Las respuestas que se apoyan en al menos un hecho en ámbar son 13 de 18 en Postmark y 12 de 20 en Mailgun, y bajan a 1 de 22 en Resend y 1 de 20 en SendGrid.

Las respuestas rotuladas entre corchetes se cuentan aparte y no se descuentan de la columna: son 17 de 21 en Amazon SES y 10 de 20 en Mailchimp.

De las cinco preguntas X, las que el instrumental no responde, seis APIs responden tres fuera de los campos.
Mailgun responde cuatro, porque es la única que publica la concurrencia que pide PI-A08, y Amazon SES dos, porque ninguna de sus fuentes publica la retención que pide PI-E03.

La versión del catálogo legible por máquina (\texttt{catalogo.yaml}) se conserva en el repositorio del trabajo junto con las 128 llamadas al motor, de modo que el procedimiento puede repetirse.
Su límite está en la clase H: las diez preguntas P se han ejecutado contra PRIME4API, pero las cuatro H cuentan como respondibles por la especificación de HARVEY4API, sin haberse ejecutado con esa herramienta.
